\documentclass[12pt,a4paper,twoside]{article}

% Configurações:

\input{conf/texto.tex}
\input{conf/matematica.tex}
\input{conf/ciencia.tex}
\input{conf/ambientes.tex}

% Para citações:

\usepackage[brazilian,hyperpageref]{backref}
\usepackage[alf]{abntex2cite}

% EDITE A PARTIR DAQUI:

\def \tituloportugues {Uso de Inteligência Artificial para detectar doenças da retina: uma aplicação usando modelos baseados em clustering e ensembles}  % Título. [EDITAR]
\def \subtitulo {Use of Artificial Intelligence to detect retinal diseases: an application using clustering-based models and ensembles}  % Título em inglês. [EDITAR]
\setcounter{page}{1} % página inicial do relatório
\def \nomeedital {PIBIC}
\def \nomerelatorio {Relatório Final de Iniciação Científica}
\def \tema {Inteligência Artificial} % pode mudar para "Eletromagnetismo", "Física Quântica" etc.
\def \coordenadas {\textit{UnB}, Brasília, 2025.} % Título em inglês. [EDITAR]
\def \paginainternet {https://proic.unb.br/}
\definecolor{cortextoquadro}{RGB}{0,58,112} % Cor do texto dos quadros do leiaute.
%\definecolor{cortextoquadro}{RGB}{0,132,61}
\definecolor{corfundoquadro}{RGB}{200,200,200} % Cor de fundo dos quadros do leiaute.
\def \fontetexto {\bfseries} % Configurações de fonte para textos dentro dos quadros do leiaute.
\def \logorevista {figs/logo-unb.png} % Defina o arquivo do logotipo.
\def \larguralogo {2.3cm} % Defina a largura com que o logotipo aparecerá na capa.

% As configurações de layout estão no arquivo mostrado abaixo:
\input{conf/layout.tex} % Não mexa nesta linha!

% Coloque abaixo os nomes dos autores e o número de suas instituições:

\author[1]{Bruno Martins Valério Bomfim}
\author[1, 2]{Profa. Dra. Marília Miranda Forte Gomes}
\author[1, 2]{Prof. Dr. Cristiano Jacques Miosso Rodrigues Mendes}

% Coloque abaixo os números e nomes das instituições:

\affil[1]{Universidade de Brasília}
\affil[2]{Programa de Pós-Graduação em Engenharia Biomédica (PPGEB/UnB)}

\input{conf/ultimos-preparativos.tex}

\begin{document}

\maketitle

\thispagestyle{capa}

% Resumo e palavras-chave. 

\hrule

\selectlanguage{portuges}
\begin{abstract}
\noindent
Este projeto de pesquisa propôs implementar e avaliar métodos de inteligência artificial (IA) para o desenvolvimento de ferramentas automatizadas destinadas ao auxílio no diagnóstico de doenças da retina, especialmente a retinopatia diabética, bem como alterações relacionadas a doenças cardiovasculares e metabólicas. A retinopatia diabética (RD) é uma complicação microvascular específica do diabetes mellitus (DM) \cite{wong_diabetic_2016, wang_diabetic_2018}, e uma das principais causas de perda de visão evitável na população adulta \cite{wong_diabetic_2016, wang_diabetic_2018, melo_current_2018, korn_malerbi_feasibility_2022, teo_global_2021}. Envolve um espectro de lesões microvasculares retinianas, mas evidências crescentes apontam também para a degeneração neuronal e processos inflamatórios como eventos precoces em sua fisiopatologia \cite{wong_diabetic_2016, wang_diabetic_2018}. Foram desenvolvidos modelos baseados em clustering, máquinas de vetores de suporte (SVM) e redes neurais convolucionais (CNN) para classificação automática das imagens retinográficas. Esses modelos foram comparados em diferentes contextos, usando tanto imagens coletadas pelo projeto quanto bases públicas. Os resultados indicam que o uso combinado dessas técnicas pode aprimorar significativamente a triagem e o diagnóstico precoce, proporcionando benefícios substanciais ao Sistema Único de Saúde (SUS), especialmente ao reduzir filas e agilizar encaminhamentos especializados. Além disso, o projeto promove avanços na engenharia biomédica e contribui para a formação de recursos humanos altamente capacitados. \\ % [EDITAR] %
\vspace{2\baselineskip}
\palavraschave{retinopatia diabética, inteligência artificial, diagnóstico automatizado, engenharia biomédica.} % [EDITAR] %
\end{abstract}
\vspace{5mm}

\hrule

\selectlanguage{portuges}

\section{Introdução} % [EDITAR]

    A diabetes mellitus (DM) tem previsão de aumento global substancial, passando de 463 milhões de indivíduos em 2019 para 700 milhões em 2045 \cite{teo_global_2021}. A retinopatia diabética (RD), complicação microvascular específica do DM \cite{wong_diabetic_2016, wang_diabetic_2018}, representa importante causa de cegueira evitável na população economicamente ativa e na terceira idade \cite{wong_diabetic_2016, wang_diabetic_2018, melo_current_2018, korn_malerbi_feasibility_2022, teo_global_2021}, com custos significativos para os indivíduos e os sistemas de saúde \cite{steinmetz_causes_2021}. Dados de 2020 estimam que 103,12 milhões de adultos convivam com RD, sendo 28,54 milhões com formas ameaçadoras à visão (VTDR). A tendência é de aumento: até 2045, esses números devem atingir 160,5 milhões e 44,82 milhões, respectivamente \cite{teo_global_2021}.

Este projeto objetiva desenvolver métodos avançados de IA para o auxílio diagnóstico dessas condições por meio da análise automatizada de imagens de retinografia, abordando tanto a identificação precoce das lesões quanto a avaliação de riscos metabólicos associados. Considerando-se que a RD apresenta fatores de risco controláveis, como hiperglicemia e hipertensão \cite{wong_diabetic_2016, melo_current_2018}, o diagnóstico precoce possibilita intervenções eficazes para impedir a progressão da doença.

A justificativa reside na capacidade da IA de reduzir custos, aumentar a acessibilidade ao diagnóstico precoce e melhorar a eficiência do atendimento no sistema público de saúde, especialmente em áreas remotas e com menor disponibilidade de especialistas. O uso de retinografia não-midriática, telemedicina e ferramentas automatizadas tem se mostrado promissor em diversos países \cite{wong_diabetic_2016, korn_malerbi_feasibility_2022, topol_high-performance_2019}, inclusive no Brasil, onde já foram testadas soluções com IA em contextos regionais como o estado de Sergipe \cite{korn_malerbi_feasibility_2022}. % [EDITAR]

\section{Referencial Teórico} % [EDITAR]

    O uso de IA em diagnóstico oftalmológico tem avançado consideravelmente nas últimas décadas, sobretudo na detecção e classificação da retinopatia diabética. \cite{gulshan_development_2016} demonstraram que algoritmos baseados em redes neurais convolucionais podem detectar essa condição com precisão comparável à avaliação humana. \cite{abramoff_pivotal_2018} ampliaram esse uso, desenvolvendo sistemas capazes de quantificar o risco de progressão da doença, facilitando intervenções médicas personalizadas.

\cite{ting_development_2017} exploraram o potencial preditivo dessas tecnologias, identificando pacientes com risco elevado de progressão para formas avançadas da retinopatia. Esses avanços demonstram a eficácia e potencial escalabilidade da IA, especialmente no rastreamento populacional e em contextos de saúde pública, como já observado em países como a Inglaterra. Atualmente, inclusive, alguns sistemas baseados em IA como o IDx foram aprovados para uso clínico autônomo em cuidados primários \cite{topol_high-performance_2019}.

No Brasil, entretanto, ainda falta uma estratégia nacional padronizada, apesar da alta prevalência de RD e do impacto significativo sobre a população adulta, conforme destacado pelas diretrizes da Sociedade Brasileira de Diabetes (SBD). O experimento realizado em Sergipe mostrou desafios como baixa adesão mesmo com triagem gratuita, necessidade de adaptação tecnológica ao contexto local e ausência de marco regulatório específico para IA em saúde \cite{korn_malerbi_feasibility_2022}. % [EDITAR]

\section{Metodologia} % [EDITAR]

    % Numeração das subseções como letras
\renewcommand{\thesubsection}{\Alph{subsection}}

Nesta seção, são descritos os métodos que orientaram as atividades desta iniciação científica. São apresentados aspectos relacionados à classificação da pesquisa, seguidos do levantamento bibliográfico e dos critérios de seleção do referencial teórico.

\subsection{Classificação da Pesquisa}
\label{ssec:classificacao-pesquisa}

Este estudo é caracterizado como uma pesquisa aplicada, de abordagem quantitativa e natureza experimental. Quanto aos objetivos, classifica-se como exploratória e explicativa, buscando compreender e explicar fenômenos relacionados à aplicação de IA no diagnóstico oftalmológico. Os procedimentos adotados incluem levantamento bibliográfico, desenvolvimento e validação de modelos computacionais, em contextos transversais e longitudinais.

\paragraph{Técnicas metodológicas empregadas}
\begin{itemize}
  \item Extração e pré-processamento de imagens retinográficas;
  \item Desenvolvimento e treinamento de modelos baseados em clustering, SVM e CNN;
  \item Avaliação comparativa de desempenho entre modelos treinados a partir de bases locais e públicas;
  \item Análise exploratória para identificar características relevantes das imagens associadas às condições oftalmológicas;
  \item Validação cruzada e testes de robustez dos modelos para garantir a aplicabilidade prática dos resultados.
\end{itemize}

\subsection{Levantamento Bibliográfico}
\label{ssec:levantamento-bibliografico}

O levantamento bibliográfico foi realizado nas bases PubMed, IEEE Xplore, Web of Science e Scopus, utilizando as seguintes strings de busca refinadas conforme a compatibilidade dos resultados com o escopo da pesquisa. A Tabela~\ref{tab:grade-completa} apresenta as consultas originais, as bases consultadas e o número de artigos obtidos, sendo filtrados apenas os resultados de publicações dos últimos dez anos (2015-2025).

\begin{table}[ht]
  \centering
  \caption{Strings de busca, bases de dados consultadas e número de resultados encontrados}
  \label{tab:grade-completa}
  \begin{tabular}{@{} l c r @{}}
    \toprule
    \textbf{String de Busca} & \textbf{Base de Dados} & \textbf{Resultados} \\ 
    \midrule
    \texttt{"diabetic retinopathy"}                     & Web of Science & 22\,528 \\ 
    \texttt{"diabetic retinopathy"}                     & Scopus         & 32\,473 \\ 
    \texttt{"diabetic retinopathy"}                     & PubMed         &  3\,542 \\ 
    \texttt{"diabetic retinopathy" AND brazil}          & Web of Science &     10 \\ 
    \texttt{"diabetic retinopathy" AND brazil}          & Scopus         &     11 \\ 
    \bottomrule
  \end{tabular}
\end{table}
\begin{flushleft}
  \small Fonte: Autor.
\end{flushleft}

\subsection{Critérios de Seleção}
\label{ssec:criterios-selecao}

Os artigos selecionados atenderam aos seguintes critérios:
\begin{itemize}
  \item Relevância científica (contribuições teóricas ou empíricas significativas);
  \item Impacto (fator de impacto dos periódicos e número de citações);
  \item Publicação recente (últimos dez anos, 2015-2025);
  \item Pertinência metodológica com o tema da pesquisa.
\end{itemize}

\section{Resultados} % [EDITAR]

    \input{texto/resultados.tex} % [EDITAR]

\section{Conclusão} % [EDITAR]

    \input{texto/conclusao.tex} % [EDITAR]

\bibliography{relatorio}

\end{document}